{%- macro tex(s) -%}
{{ s | replace('\\', '\\textbackslash{}') | replace('&', '\\&') | replace('%', '\\%') | replace('$', '\\$') | replace('#', '\\#') | replace('_', '\\_') }}
{%- endmacro -%}
{%- macro url(s) -%}
{{ s | replace('%', '\\%') | replace('#', '\\#') }}
{%- endmacro -%}
{%- set paper = cookiecutter.paper_size -%}
{% raw %}%-----------------------------------------------------------------------------------------------------------------------------------------------%
%	The MIT License (MIT)
%
%	Copyright (c) 2021 Philip Empl
%
%	Permission is hereby granted, free of charge, to any person obtaining a copy
%	of this software and associated documentation files (the "Software"), to deal
%	in the Software without restriction, including without limitation the rights
%	to use, copy, modify, merge, publish, distribute, sublicense, and/or sell
%	copies of the Software, and to permit persons to whom the Software is
%	furnished to do so, subject to the following conditions:
%
%	The above copyright notice and this permission notice shall be included in
%	all copies or substantial portions of the Software.
%
%	THE SOFTWARE IS PROVIDED "AS IS", WITHOUT WARRANTY OF ANY KIND, EXPRESS OR
%	IMPLIED, INCLUDING BUT NOT LIMITED TO THE WARRANTIES OF MERCHANTABILITY,
%	FITNESS FOR A PARTICULAR PURPOSE AND NONINFRINGEMENT. IN NO EVENT SHALL THE
%	AUTHORS OR COPYRIGHT HOLDERS BE LIABLE FOR ANY CLAIM, DAMAGES OR OTHER
%	LIABILITY, WHETHER IN AN ACTION OF CONTRACT, TORT OR OTHERWISE, ARISING FROM,
%	OUT OF OR IN CONNECTION WITH THE SOFTWARE OR THE USE OR OTHER DEALINGS IN
%	THE SOFTWARE.
%
%-----------------------------------------------------------------------------------------------------------------------------------------------%
%
{% endraw %}%	{{ tex(cookiecutter.full_name) }} -- CV
{% raw %}%	Based on the "sidebar" modern-latex-cv template by Philip Empl
%	(https://www.overleaf.com/latex/templates/modern-latex-cv/qmdwjvcrcrph)
%
%============================================================================%
%
%	DOCUMENT DEFINITION
%
%============================================================================%

{% endraw %}\documentclass[10pt,{{ paper }}paper,english]{article}
{% raw %}

%----------------------------------------------------------------------------------------
%	ENCODING
%----------------------------------------------------------------------------------------

\usepackage[utf8]{inputenc}
\usepackage[USenglish]{isodate}
\usepackage{fancyhdr}
\usepackage[numbers]{natbib}

%----------------------------------------------------------------------------------------
%	LOGIC
%----------------------------------------------------------------------------------------

\usepackage{xstring, xifthen}
\usepackage{enumitem}
\usepackage[english]{babel}
\usepackage{pdfpages}
\usepackage{changepage}

%----------------------------------------------------------------------------------------
%	FONT BASICS
%----------------------------------------------------------------------------------------

\usepackage[default]{raleway}

\renewcommand*\familydefault{\sfdefault}
\usepackage[T1]{fontenc}

\usepackage{moresize}

%----------------------------------------------------------------------------------------
%	FONT AWESOME ICONS
%----------------------------------------------------------------------------------------

\usepackage{fontawesome}

\newcommand{\vcenteredinclude}[1]{\begingroup
\setbox0=\hbox{\includegraphics{#1}}%
\parbox{\wd0}{\box0}\endgroup}
\newcommand{\tab}[1]{\hspace{.2\textwidth}\rlap{#1}}
\newcommand*{\vcenteredhbox}[1]{\begingroup
\setbox0=\hbox{#1}\parbox{\wd0}{\box0}\endgroup}

\newcommand{\icon}[3] {
	\makebox(#2, #2){\textcolor{maincol}{\ifthenelse{\equal{#1}{Linkedin}}{\raisebox{3.5pt}{\csname fa#1\endcsname}}{\csname fa#1\endcsname}}}
}

\newcommand{\icontext}[4] {
	\vcenteredhbox{\icon{#1}{#2}{#3}}  \hspace{2pt}  \parbox{0.9\mpwidth}{\textcolor{#4}{#3}}
}

\newcommand{\iconhref}[5] {
    \vcenteredhbox{\icon{#1}{#2}{#5}}  \hspace{2pt} \href{#4}{\textcolor{#5}{#3}}
}

\newcommand{\iconemail}[5] {
    \vcenteredhbox{\icon{#1}{#2}{#5}}  \hspace{2pt} \href{mailto:#4}{\textcolor{#5}{#3}}
}

%----------------------------------------------------------------------------------------
%	PAGE LAYOUT  DEFINITIONS
%----------------------------------------------------------------------------------------

\usepackage{paracol}
\usepackage{tikzpagenodes}
\usetikzlibrary{calc}
\usepackage{lmodern}
\usepackage{multicol}
\usepackage{atbegshi}
{% endraw %}\usepackage[{{ paper }}paper]{geometry}
{% raw %}
\geometry{top=1cm, bottom=0.5cm, left=1cm, right=1cm}

\usepackage{fancyhdr}
\pagestyle{empty}

\setlength{\parindent}{0mm}
\renewcommand{\baselinestretch}{0.94}

%----------------------------------------------------------------------------------------
%	TABLE /ARRAY DEFINITIONS
%----------------------------------------------------------------------------------------

\usepackage{array}

\newcolumntype{x}[1]{%
>{\raggedleft\hspace{0pt}}p{#1}}%

%----------------------------------------------------------------------------------------
%	GRAPHICS DEFINITIONS
%----------------------------------------------------------------------------------------

\usepackage{graphicx}
\usepackage{tikz}
\usepackage{ragged2e}
\usetikzlibrary{shapes, backgrounds,mindmap, trees}
\graphicspath{{src/resources/}{resources/}}

%----------------------------------------------------------------------------------------
%	Color DEFINITIONS
%----------------------------------------------------------------------------------------
\usepackage{transparent}
\usepackage{color}

\definecolor{maincol}{RGB}{ 64,64,64}
\definecolor{darkcol}{RGB}{ 70, 70, 70 }
\definecolor{lightcol}{RGB}{245,245,245}
\definecolor{accentcol}{RGB}{59,77,97}

\usepackage[hidelinks]{hyperref}

\newcommand{\mpwidth}{\linewidth-\fboxsep-\fboxsep}

%============================================================================%
%
%	CV COMMANDS
%
%============================================================================%

\newcommand{\cvlist}[1] {
	\begin{itemize}{#1}\end{itemize}
}

\newcommand{\cvtext}[1] {
	\begin{tabular*}{1\mpwidth}{p{0.98\mpwidth}}
		\parbox{1\mpwidth}{\RaggedRight #1}
	\end{tabular*}
}
\newcommand{\cvtextsmall}[1] {
	\begin{tabular*}{0.8\mpwidth}{p{0.8\mpwidth}}
		\parbox{0.8\mpwidth}{\RaggedRight #1}
	\end{tabular*}
}

\newcommand{\cvsection}[1] {
	\vspace{10pt}
	\cvtext{
		\textbf{\LARGE{\textcolor{darkcol}{#1}}}\\[-4pt]
		\textcolor{accentcol}{ \rule{0.2\textwidth}{1.5pt} } \\
	}
}

\newcommand{\cvsectionsmall}[1] {
	\vspace{10pt}
	\cvtext{
		\textbf{\Large{\textcolor{darkcol}{#1}}}\\[-4pt]
		\textcolor{accentcol}{ \rule{0.2\textwidth}{1.5pt} } \\
	}
}

\newcommand{\cvsidesection}[1] {
	\vspace{8pt}
	\cvtext{
		\textbf{\large{\textcolor{darkcol}{#1}}}\\[-6pt]
		\textcolor{accentcol}{ \rule{0.2\textwidth}{1.2pt} } \\[4pt]
	}
}

\newcommand{\cvheadline}[1] {
	\vspace{16pt}
	\cvtext{
		\textbf{\Huge{\textcolor{accentcol}{#1}}}\\[-4pt]

	}
}

\newcommand{\cvsubheadline}[1] {
	\vspace{16pt}
	\cvtext{
		\textbf{\huge{\textcolor{darkcol}{#1}}}\\[-4pt]

	}
}

%----------------------------------------------------------------------------------------
%	META SKILL
%----------------------------------------------------------------------------------------

\newcommand{\cvskill}[2] {
	\begin{tabular*}{1\mpwidth}{p{0.72\mpwidth}  r}
 		\textcolor{black}{\textbf{#1}} & \textcolor{maincol}{#2}\\
	\end{tabular*}%
}


%----------------------------------------------------------------------------------------
%	 CV EVENT
%----------------------------------------------------------------------------------------

\newcommand{\cvevent}[4] {

	\parbox{\mpwidth}{
		\begin{tabular*}{1\mpwidth}{p{0.66\mpwidth}  r}
	 		\textcolor{black}{\textbf{#2}} & \colorbox{accentcol}{\makebox[0.32\mpwidth]{\textcolor{white}{\textbf{#1}}}} \\
			\textcolor{maincol}{#3} & \\
		\end{tabular*}\\[3pt]

		\ifthenelse{\isempty{#4}}{}{
			\cvtext{#4}\\
		}
	}
	\vspace{4pt}
}

%----------------------------------------------------------------------------------------
%	 CV META EVENT
%----------------------------------------------------------------------------------------

\newcommand{\cvmetaevent}[4] {
	\textcolor{maincol} { \cvtext{\textbf{#1}}}\\%
	\ifthenelse{\isempty{#2}}{}{\textcolor{black} {\cvtext{\textbf{#2}} }\\}%
	\ifthenelse{\isempty{#3}}{}{\cvtext{{ \textcolor{maincol} {#3} }}}%
	\ifthenelse{\isempty{#4}}{}{\\\cvtext{#4}}%
	\par\vspace{10pt}%
}

% HEADER AND FOOTER
%====================================
\newcommand\Header[1]{%
\begin{tikzpicture}[remember picture,overlay]
\fill[accentcol]
  (current page.north west) -- (current page.north east) --
  ([yshift=50pt]current page.north east|-current page text area.north east) --
  ([yshift=50pt,xshift=-3cm]current page.north|-current page text area.north) --
  ([yshift=10pt,xshift=-5cm]current page.north|-current page text area.north) --
  ([yshift=10pt]current page.north west|-current page text area.north west) -- cycle;
\node[font=\sffamily\bfseries\color{white},anchor=west,
  xshift=0.7cm,yshift=-0.32cm] at (current page.north west)
  {\fontsize{12}{12}\selectfont {#1}};
\end{tikzpicture}%
}

\newcommand\Footer[1]{%
\begin{tikzpicture}[remember picture,overlay]
\fill[lightcol]
  (current page.south east) -- (current page.south west) --
  ([yshift=-80pt]current page.south east|-current page text area.south east) --
  ([yshift=-80pt,xshift=-6cm]current page.south|-current page text area.south) --
  ([xshift=-2.5cm,yshift=-10pt]current page.south|-current page text area.south) --
  ([yshift=-10pt]current page.south east|-current page text area.south east) -- cycle;
\end{tikzpicture}%
}

%============================================================================%
{% endraw %}%
%	DOCUMENT CONTENT
%
%  !! EDIT YOUR CV BELOW THIS LINE !!
%  Everything above this block is template boilerplate. Only touch it if you
%  want to change the design (colours live in "Color DEFINITIONS", margins in
%  "PAGE LAYOUT DEFINITIONS").
%
%============================================================================%
\begin{document}

\columnratio{0.27}
\setlength{\columnsep}{1.6em}
\setlength{\columnseprule}{1pt}
\colseprulecolor{white}

\AtBeginShipoutFirst{\Header{Curriculum Vitae}\Footer{}}
\AtBeginShipout{\AtBeginShipoutAddToBox{\Header{Curriculum Vitae}\Footer{}}}

\newpage

\colseprulecolor{lightcol}
\columnratio{0.27}
\setlength{\columnsep}{1.6em}
\setlength{\columnseprule}{1pt}
\begin{paracol}{2}

\begin{leftcolumn}
{%- if cookiecutter.include_photo %}
%---------------------------------------------------------------------------------------
%	PHOTO (cropped to a circle, bordered)
%   Drop your picture at src/resources/photo.jpg (a square-ish crop works best,
%   the template circles it automatically). Until that file exists, a circle
%   with your initials is shown instead.
%----------------------------------------------------------------------------------------
\newlength{\photodiam}
\setlength{\photodiam}{0.85\linewidth}
\newif\ifhavephoto
\IfFileExists{src/resources/photo.jpg}{\havephototrue}{\IfFileExists{resources/photo.jpg}{\havephototrue}{\havephotofalse}}
\begin{center}
\begin{tikzpicture}
	\ifhavephoto
	\begin{scope}
		\clip (0,0) circle (0.5\photodiam);
		\node[inner sep=0pt] at (0,0) {\includegraphics[height=\photodiam,keepaspectratio]{photo.jpg}};
	\end{scope}
	\else
	\fill[lightcol] (0,0) circle (0.5\photodiam);
	\node[text=accentcol, font=\Huge\bfseries] at (0,0) {{ '{' }}{% for part in cookiecutter.full_name.split() %}{{ part[0] }}{% endfor %}{{ '}' }};
	\fi
	\draw[color=accentcol!45!white, line width=2.5pt] (0,0) circle (0.5\photodiam);
\end{tikzpicture}
\end{center}
{%- else %}
\vspace{4pt}
{%- endif %}

%---------------------------------------------------------------------------------------
%	NAME / TITLE
%   Edit the text inside { } on each line.
%   The title is split across two lines so it doesn't wrap awkwardly.
%   \\[3pt] is the gap between the lines; increase the number for more space.
%----------------------------------------------------------------------------------------
	\fcolorbox{white}{white}{\begin{minipage}[c][1.7cm][c]{1\mpwidth}
		\LARGE{\textbf{\textcolor{maincol}{ {{- tex(cookiecutter.full_name) -}} }}} \\[2pt]
		\normalsize{ \textcolor{maincol} { {{- tex(cookiecutter.headline) -}} } }
{%- if cookiecutter.headline_secondary %} \\[3pt]
		\normalsize{ \textcolor{maincol} { {{- tex(cookiecutter.headline_secondary) -}} } }
{%- endif %}
\end{minipage}} \\

\cvsidesection{Contact}

% Each line is: icon, size (leave as 14), display text, colour (leave as black).
% \iconemail: the 4th argument is the actual address used for the mailto: link.
% \iconhref:  the 4th argument is the URL; edit both the URL and the display text.
% Icon names come from the fontawesome package (v4): MapMarker, Phone, Envelope,
% Linkedin, Github, Globe, ... (drop the \fa prefix).
\icontext{MapMarker}{14}{ {{- tex(cookiecutter.location) -}} }{black}\\[3pt]
{%- if cookiecutter.phone %}
\icontext{Phone}{14}{ {{- tex(cookiecutter.phone) -}} }{black}\\[3pt]
{%- endif %}
\iconemail{Envelope}{14}{ {{- tex(cookiecutter.email) -}} }{ {{- url(cookiecutter.email) -}} }{black}\\[3pt]
{%- if cookiecutter.linkedin_url %}
\iconhref{Linkedin}{14}{ {{- tex(cookiecutter.full_name) -}} }{ {{- url(cookiecutter.linkedin_url) -}} }{black}\\[3pt]
{%- endif %}
{%- if cookiecutter.github_username %}
\iconhref{Github}{14}{ {{- tex(cookiecutter.github_username) -}} }{https://github.com/ {{- cookiecutter.github_username -}} }{black}\\[3pt]
{%- endif %}

\cvsidesection{Skills}

% Each \item is one bullet. Add, remove, or edit lines freely.
% \textbf{...} makes text bold: use it for the keywords recruiters scan for.
% \mbox{...} keeps a phrase on one line so it doesn't break awkwardly.
\cvtext{
\begin{itemize}[leftmargin=*,topsep=2pt,itemsep=2pt,parsep=0pt]
	\item \textbf{Python}, \textbf{Go}, \textbf{SQL}
	\item \textbf{Kubernetes}, \textbf{Terraform}, \mbox{\textbf{GitHub Actions}}
	\item \textbf{Linux} administration \& \mbox{troubleshooting}
	\item Technical writing \& documentation
	\item Fluent \textbf{English} (C1), native Polish
\end{itemize}
}

%---------------------------------------------------------------------------------------
%	INTERESTS
%   Simple list: each \item is one line.
%----------------------------------------------------------------------------------------
\cvsidesection{Interests}

\cvtext{
\begin{itemize}[leftmargin=*,topsep=0pt,itemsep=0pt,parsep=0pt]
	\item Open source
	\item Distributed systems
	\item Mountain hiking
	\item Photography
\end{itemize}
}

%---------------------------------------------------------------------------------------
%	EDUCATION
%   Each entry uses \cvmetaevent with four lines:
%     {years}
%     {Degree or qualification}
%     {Institution}
%     {Subtitle: wrap in \textit{...} for italics, or leave empty {}}
%   To add another entry, copy a full \cvmetaevent block and edit the four lines.
%----------------------------------------------------------------------------------------
\cvsidesection{Education}

\cvmetaevent
{2019 -- 2021}
{MSc, Computer Science}
{University of Example}
{\textit{Distributed Systems}}

\vspace{2pt}
\cvmetaevent
{2015 -- 2019}
{BSc, Computer Science}
{University of Example}
{}

\end{leftcolumn}
\begin{rightcolumn}
%---------------------------------------------------------------------------------------
%	PROFILE
%   One-paragraph summary, the first thing a recruiter reads.
%   Keep it to 3-4 sentences and tailor it to the role you're applying for.
%----------------------------------------------------------------------------------------
\cvsection{Profile}
\vspace{2pt}

\cvtext{Software engineer with six years of experience building reliable cloud infrastructure and developer tooling. Comfortable owning a system end to end, from design and implementation to on-call and documentation. Looking for a team where \mbox{platform engineering} directly shapes how quickly product teams can ship.}

%---------------------------------------------------------------------------------------
%	WORK EXPERIENCE
%   Each job uses \cvevent with four parts, in this order:
%     {date range}    shown in the dark badge on the right
%     {Job title}
%     {Company name}
%     {Description}   either \cvtext{one paragraph} or an itemize list (see below)
%
%   Add a job: copy a full \cvevent block. Remove a job: delete the block.
%   Keep jobs in reverse chronological order (newest first).
%   Too many bullets will push the CV onto a second page.
%
%   Special characters: \& = &, \% = %, \L{} = Ł, \'c = ć, \k{e} = ę, etc.
%   (UTF-8 also works directly: Ł, ć, ę, ń, ó, ś, ź, ż.)
%----------------------------------------------------------------------------------------

\vspace{4pt}
\cvsection{Work experience}
\vspace{2pt}

\cvevent
{03/2023 -- Present}
	{Senior Software Engineer}
	{Example Corp}
	{
	\begin{itemize}[leftmargin=*,topsep=0pt,itemsep=0pt,parsep=0pt]
		\item Led the migration of 40+ services to \textbf{Kubernetes}, cutting deployment time from hours to minutes.
		\item Designed a self-service \textbf{Terraform} module catalogue used by every product team.
		\item Introduced SLOs and on-call runbooks, reducing pages by 60\% in the first quarter.
		\item Mentored three engineers through their first production launches.
	\end{itemize}
	}

\cvevent
{07/2021 -- 02/2023}
	{Software Engineer}
	{Sample Systems}
	{
	\begin{itemize}[leftmargin=*,topsep=0pt,itemsep=0pt,parsep=0pt]
		\item Built internal CLI tooling in \textbf{Go} that automated environment provisioning.
		\item Maintained CI/CD pipelines in \textbf{GitHub Actions} for a monorepo of 25 services.
	\end{itemize}
	}

\cvevent
{10/2019 -- 06/2021}
	{Software Engineering Intern}
	{Placeholder Labs}
	{
	% A short role: \cvtext{} gives a single paragraph instead of bullets.
	\cvtext{Wrote data pipelines in \textbf{Python} and \textbf{SQL} that replaced a weekly manual reporting process.}
	}

%---------------------------------------------------------------------------------------
%	PROJECTS
%   \cvevent works for anything with a title, an organisation and dates.
%   Rename the section (e.g. Publications, Certifications) or delete it.
%----------------------------------------------------------------------------------------
\vspace{4pt}
\cvsection{Projects}
\vspace{2pt}

\cvevent
{2020 -- Present}
	{Maintainer, example-tool}
	{Open source, github.com/example/example-tool}
	{
	\cvtext{A small command line tool for validating configuration files, with 1k+ stars and contributors from a dozen companies.}
	}

\cvevent
{2022}
	{Speaker, ``Shipping Kubernetes Upgrades Without Pager Duty''}
	{Example Cloud Native Meetup}
	{
	\cvtext{A 30 minute talk on canary node pools, automated rollbacks and the runbooks that made weekly upgrades boring.}
	}

\cvsection{Certifications}
\vspace{2pt}

\cvevent
{2023}
	{Certified Kubernetes Administrator (CKA)}
	{The Linux Foundation}
	{}

\cvevent
{2022}
	{AWS Certified Solutions Architect, Associate}
	{Amazon Web Services}
	{}
{%- if cookiecutter.include_gdpr_clause %}

% Polish GDPR (RODO) recruitment consent clause, printed in tiny text at the
% bottom of the page. Many Polish employers require it; delete the block if
% you apply abroad.
{% raw %}\begin{tikzpicture}[remember picture,overlay]
\node[anchor=south west, xshift=1cm, yshift=0.3cm, text width=\paperwidth-2cm, align=left,
  font=\fontsize{4}{3.2}\selectfont, text=maincol!45] at (current page.south west) {%
Wyra\.zam zgod\k{e} na przetwarzanie moich danych osobowych dla potrzeb niezb\k{e}dnych do realizacji procesu tej oraz przysz\l{}ych rekrutacji (zgodnie z ustaw\k{a} z dnia 10 maja 2018 roku o ochronie danych osobowych (Dz.\ Ustaw z 2018, poz.\ 1000) oraz zgodnie z Rozporz\k{a}dzeniem Parlamentu Europejskiego i Rady (UE) 2016/679 z dnia 27 kwietnia 2016 r.\ w sprawie ochrony os\'ob fizycznych w zwi\k{a}zku z przetwarzaniem danych osobowych i w sprawie swobodnego przep\l{}ywu takich danych oraz uchylenia dyrektywy 95/46/WE (RODO)).%
};
\end{tikzpicture}
{% endraw %}{%- endif %}

\end{rightcolumn}
\end{paracol}

\end{document}
